\chapter*{Abstract}
\addcontentsline{toc}{chapter}{Abstract}

This thesis investigates numerical methods for solving nonlinear equations and systems of nonlinear equations. Such problems arise frequently in various areas of science, engineering, optimization, and computational mathematics, and in many practical cases their exact analytical solutions cannot be obtained. For this reason, numerical approaches play an important role in approximating their solutions.

The study focuses on both single nonlinear equations and systems of nonlinear equations. For single-variable problems, the bisection method, secant method, Newton's method, damped Newton method, and Brent's method are examined. For nonlinear systems, multivariable Newton and Broyden methods are studied. The theoretical foundations, convergence properties, advantages, limitations, and failure cases of these methods are discussed.

In addition to the theoretical analysis, the selected numerical methods are implemented in Python using a modular computational structure. Numerical experiments are performed in order to compare the methods in terms of convergence behavior, residual reduction, and stability. Convergence plots, residual analyses, and iteration histories are used to evaluate the performance of the algorithms.

The thesis also investigates the relationship between nonlinear equation solving and optimization techniques. In particular, residual minimization, Newton-type optimization methods, damping strategies, line-search techniques, and the Armijo condition are discussed from an optimization perspective. An adaptive damped Newton method based on Armijo backtracking is also implemented and analyzed numerically.

Overall, the study combines theoretical numerical analysis with practical computational implementation and provides a structured examination of numerical methods for nonlinear problems.